\section{Introduction}
\label{sec:introduction}

\PARstart{E}{dge-cloud} Internet of Things (IoT) systems give a scheduler two imperfect execution sites. Edge nodes sit close to devices and can avoid long communication paths, but their CPU, memory, queue, and energy budgets are limited. Cloud nodes provide larger compute pools, yet cloud execution depends on network delay, uplink capacity, packet loss, jitter, cold starts, and maintenance state. The scheduling decision is not a fixed preference for edge or cloud; it is repeated under changing resource and channel conditions.

The problem becomes harder when arrivals are bursty. A camera-heavy interval, a firmware rollout, or an emergency burst can make many devices compete for the same edge resources at nearly the same time. A latency-only policy can look strong on average and still fail in these windows. Routing too much work to edge may reduce communication delay while increasing queue delay and admission-failure risk. Routing too much work to cloud may avoid local overload but miss tight deadlines when the network is degraded. A usable policy has to learn these tradeoffs without assuming that every zone sees the same workload distribution.

Deep reinforcement learning (DRL) is a natural fit for this setting because task offloading is sequential and state dependent. Double Deep Q-Networks (DDQNs) are especially relevant when the action space is discrete, because the double target reduces overestimation compared with a single Q-network. A centralized DDQN, however, assumes that all training data can be pooled and that one global distribution is enough. Edge-cloud IoT systems often violate both assumptions. Different edge zones can see different task types, channel profiles, device classes, and failure patterns. Federated learning offers a better training structure: each zone can train locally on its own distribution, and the server aggregates model updates rather than raw data.

This work studies task offloading as a federated DDQN problem. The Stage-1 model is a zone-aware \fedddqn{} policy for binary edge-cloud offloading. The action space is deliberately small: \EdgeAction{} selects edge execution and \CloudAction{} selects cloud execution. Rejection is not a third action. It is a derived admission or infeasibility outcome that occurs when the selected path cannot satisfy resource, deadline, security, or state constraints. This distinction matters. If rejection were an action, a scheduler could appear cleaner by refusing difficult tasks. In the binary setting, the policy must decide where feasible work should go.

The paper also includes a second-stage resource allocation model. After \fedddqn{} chooses edge execution, a rejection-aware residual allocator assigns CPU, memory, and bandwidth shares. The allocator starts from a demand-proportional allocation, learns deadline, priority, and QoS/admission-risk corrections, and then applies capacity normalization. It does not replace the offloading policy. It manages resources for tasks that the proposed scheduler has already selected for local execution.

The evaluation uses \dataset{}, a realistic synthetic benchmark built for controlled comparison. The benchmark is not measured deployment telemetry. It is more stressful than independent random task rows because it includes workload phases, temporal bursts, outages, jitter storms, congestion, edge degradation, cloud maintenance, and queue carryover. The fixed dataset keeps comparisons repeatable while still exposing policies to deployment-like variation.

On the untouched temporal test split, \fedddqn{} achieved \FedLatency{} ms average latency, \FedSLA\% SLA satisfaction, \FedRejection\% rejection, and \FedEdgeUse\% edge usage. Its latency is \FedVsDDQN\% lower than centralized DDQN and \FedVsFLDDPG\% lower than the FL-DDPG-style baseline. The proposed model gives the lowest latency among the main comparison methods reported in Table~\ref{tab:main_results} while preserving a federated, zone-aware training design.

Three questions guide the study. First, can federated DDQN reduce test latency compared with centralized DDQN and the implemented offloading baselines when the workload is evaluated in a future temporal interval? Second, can a reward that penalizes SLA miss and admission failure avoid the failure mode in which mean latency improves only because difficult tasks become infeasible? Third, once the destination is selected, can a residual allocator improve CPU, memory, and bandwidth assignment without changing that destination? These questions determine the paper structure: the method separates offloading from allocation, and the results report latency, QoS, rejection, scenario behavior, ablation, and allocation metrics rather than a single latency number.

The contributions are as follows:
\begin{itemize}
\item A zone-aware \fedddqn{} framework for binary edge-cloud IoT task offloading under non-identical edge-zone distributions.
\item A channel-aware and rejection-aware reward design that accounts for latency, SLA misses, infeasibility risk, edge overuse, and heavy-task misrouting.
\item A fixed-schema realistic synthetic benchmark, \dataset{}, with temporal bursts, contention effects, network faults, edge degradation, and cloud maintenance.
\item A two-stage decision pipeline in which \fedddqn{} selects edge or cloud execution and a residual allocator refines CPU, memory, and bandwidth shares for edge-selected tasks.
\item A test-only evaluation against a faithful DDQN baseline, a federated actor--critic-style comparator, several style-inspired scheduling and metaheuristic comparators, and a non-trainable latency-reference Oracle, with scenario-wise, multi-seed, ablation, and allocator analyses.
\end{itemize}

The manuscript proceeds as follows. Section~\ref{sec:related_work} reviews edge-cloud offloading, DRL, federated learning, and resource allocation. Section~\ref{sec:system_model} defines the system model and optimization objective. Section~\ref{sec:dataset} describes \dataset{} and the feature pipeline. Section~\ref{sec:fed_ddqn} presents the proposed federated DDQN offloading method. Section~\ref{sec:stage2_allocator} describes the resource allocator. Section~\ref{sec:experimental_setup} gives the experimental protocol, and Section~\ref{sec:results} reports the results and discussion.
